\documentclass[12pt, a4paper]{article}
\usepackage[utf8]{inputenc}
\usepackage[spanish]{babel}
\usepackage{graphicx}
\usepackage{hyperref}
\usepackage{geometry}
\usepackage{listings}
\usepackage{xcolor}

\geometry{a4paper, margin=1in}

\definecolor{codegray}{rgb}{0.5,0.5,0.5}
\definecolor{backcolour}{rgb}{0.95,0.95,0.92}
\lstdefinestyle{mystyle}{
    backgroundcolor=\color{backcolour},   
    commentstyle=\color{codegray},
    basicstyle=\ttfamily\footnotesize,
    breakatwhitespace=false,         
    breaklines=true,                 
    captionpos=b,                    
    keepspaces=true,                 
    numbers=left,                    
    numbersep=5pt,                  
    showspaces=false,                
    showstringspaces=false,
    showtabs=false,                  
    tabsize=2
}
\lstset{style=mystyle}

\title{Documento de Investigación Técnica y Manual Operativo:\\
\large Implementación y Uso de Pull Requests (PR) en Repositorios Colaborativos}
\author{Vicente Matus \and Daniela Romero \and Renato Carrasco \and \\ Marcelo Matamala \and Fabian Sánchez \and Esban Vejar}
\date{\today}

\begin{document}

\maketitle

\begin{abstract}
El presente documento técnico aborda la implementación metodológica y el uso práctico de los \textit{Pull Requests} (PR) como mecanismo de integración continua y control de calidad de código. Dirigido al equipo de desarrollo y al administrador del repositorio del Taller de Integración III, el documento establece de manera exhaustiva y paso a paso cómo configurar políticas de protección de ramas en la interfaz de GitHub, evitando inyecciones directas de código. Posteriormente, actúa como un manual de usuario detallando el ciclo de vida de un PR: apertura de solicitudes, revisiones cruzadas (Code Review), metodologías de integración (Squash and Merge) y resolución de conflictos estructurales.
\end{abstract}

\tableofcontents
\newpage

\section{Introducción: La Filosofía del Pull Request}
A diferencia de los comandos tradicionales de Git (\texttt{commit}, \texttt{push}, \texttt{pull}), un \textbf{Pull Request (PR)} no pertenece a la tecnología base de Git, sino que es una característica de plataformas de alojamiento como GitHub. Estructuralmente, es una "petición formal" que un desarrollador emite para que el código contenido en su rama aislada sea inspeccionado y, si cumple los estándares, sea integrado a las ramas de despliegue (\texttt{develop} o \texttt{main}).

En el contexto de la Ingeniería de Software, el PR no es un mero trámite administrativo, sino un \textbf{punto crítico de control de calidad}. Garantiza que ningún código ingrese a la base principal sin haber sido auditado al menos por un segundo par de ojos.

\section{Fase de Implementación: Configuración del Encargado (Administrador)}
Para que el flujo de Pull Requests tenga rigor metodológico, es mandatario deshabilitar técnicamente la capacidad de alterar las ramas principales de forma unilateral. Este procedimiento debe ser ejecutado \textbf{exclusivamente por el administrador del repositorio}.

Instrucciones exactas en la interfaz de GitHub:
\begin{enumerate}
    \item Ingresar a la interfaz gráfica del repositorio y hacer clic en la pestaña superior \textbf{Settings} (ícono de rueda dentada).
    \item En el menú de navegación lateral izquierdo, bajo la categoría \textit{Code and automation}, seleccionar \textbf{Branches}.
    \item Ejecutar la acción haciendo clic en el botón derecho \textbf{Add branch protection rule}.
    \item En el campo de texto \textbf{Branch name pattern}, escribir el término exacto: \texttt{develop}.
    \item En el listado inferior, marcar la directiva \textbf{Require a pull request before merging}. Su activación desplegará parámetros secundarios.
    \item Asegurar que la sub-casilla \textbf{Require approvals} esté marcada, y verificar que el contador numérico de aprobaciones esté seteado en \textbf{1}.
    \item (Sugerido) Marcar la casilla \textbf{Require conversation resolution before merging} para bloquear integraciones si existen comentarios técnicos pendientes de responder.
    \item Desplazarse al pie de la página y confirmar con el botón \textbf{Create} o \textbf{Save changes}.
    \item Repetir exhaustivamente los pasos 3 al 8 sustituyendo el patrón de texto por \texttt{main}, protegiendo así la rama de producción.
\end{enumerate}

\textbf{Efecto Técnico:} A partir de este hito de configuración, si un integrante ejecuta localmente \texttt{git push origin develop}, el servidor rechazará la petición (\textit{Push rejected}), forzando la apertura obligatoria de un Pull Request.

\section{Manual Operativo: El Ciclo de Vida del Pull Request}
El siguiente flujo instruye el actuar del desarrollador base y del revisor asignado. Se asume que el autor ha ejecutado un \texttt{git push origin tu-rama} en su entorno local.

\subsection{Paso A: Apertura del PR (Rol: Autor)}
\begin{itemize}
    \item Al acceder a GitHub.com, la plataforma presentará un banner superior de color amarillo notificando la subida de la rama. Hacer clic en el botón \textbf{Compare \& pull request}.
    \item En el encabezado superior de la nueva interfaz, verificar el flujo de integración en los menús desplegables: la variable \textbf{base} debe apuntar a \texttt{develop}, mientras que la variable \textbf{compare} debe reflejar la rama de trabajo aislada.
    \item \textbf{Trazabilidad Documental:} Es obligatorio redactar el título y la descripción del PR de forma descriptiva, indicando la casuística resuelta y qué parámetros de prueba debe utilizar el revisor.
    \item En la barra lateral derecha, en la sección \textbf{Reviewers}, se debe seleccionar el nombre de usuario del compañero o encargado responsable de auditar el código.
    \item Finalizar presionando el botón \textbf{Create pull request}.
\end{itemize}

\subsection{Paso B: Revisión de Código (Rol: Revisor o Encargado)}
Este proceso demanda rigor analítico por parte del auditor.
\begin{itemize}
    \item El revisor ingresa al PR y selecciona la pestaña \textbf{Files changed}. Esta vista contrastará el código obsoleto (en rojo) frente a las adiciones (en verde).
    \item Durante el escrutinio, el revisor debe posicionar el cursor sobre el número de una línea de código anómala; emergerá un icono azul con un signo más (\textbf{+}).
    \item Al hacer clic en el \textbf{+}, se despliega una caja de texto para documentar el error encontrado. Confirmar presionando \textbf{Start a review}.
    \item Una vez finalizado el barrido del documento, el revisor debe hacer clic en el botón superior derecho \textbf{Review changes}, seleccionando la sentencia correspondiente:
    \begin{itemize}
        \item \textit{Comment:} Observaciones neutras que no detienen el flujo.
        \item \textit{Approve:} Conformidad técnica absoluta. Habilita el botón de fusión.
        \item \textit{Request changes:} Veto técnico temporal. Bloquea el botón de fusión hasta que el autor rectifique los errores reportados.
    \end{itemize}
    \item Confirmar con \textbf{Submit review}.
\end{itemize}

\subsection{Paso C: Rectificación (Rol: Autor)}
Si el PR fue objetado (\textit{Request changes}), el autor \textbf{no debe clausurar el Pull Request ni generar uno nuevo}.
Deberá retornar a su entorno local (ej. VSCode), reparar la lógica en la misma rama afectada, generar un \texttt{git commit -m "fix: correcciones"} e inyectarlo vía \texttt{git push}. El portal de GitHub actualizará automáticamente la visualización del PR para una re-evaluación.

\subsection{Paso D: Fusión y Despliegue (Rol: Encargado / Autor)}
Con el dictamen favorable (\textit{Approved}), la interfaz habilitará las opciones de integración.
\begin{itemize}
    \item Desplazarse al final de la conversación del PR. Junto al botón verde de "Merge pull request", hacer clic en la flecha de menú desplegable.
    \item Seleccionar la modalidad \textbf{Squash and merge}.
    \item \textit{Justificación Arquitectónica:} Esta estrategia algorítmica toma el historial granular de la rama ("intento 1", "bugfix") y lo condensa matemáticamente en un único \textit{commit} profesional y semántico al inyectarlo a \texttt{develop}.
    \item Presionar \textbf{Confirm squash and merge}.
    \item Acto seguido, la interfaz sugerirá un botón \textbf{Delete branch}. Presionarlo es imperativo para mantener la salubridad del repositorio, borrando la rama remota que ya cumplió su vida útil.
\end{itemize}

\section{Anexo: Resolución de Conflictos Estructurales}
Si la interfaz de GitHub advierte \textbf{"Can't automatically merge"}, significa que el autor y un agente externo editaron concurrentemente las mismas líneas lógicas de un archivo idéntico.

El protocolo local de desambiguación consta de 5 pasos:
\begin{enumerate}
    \item Asegurar la posición en la rama problemática: \texttt{git checkout tu-rama}.
    \item Importar las novedades desde el servidor: \texttt{git pull origin develop}.
    \item El terminal de Git abortará la fusión. Al abrir el editor (VSCode), el código en disputa estará confinado entre marcadores \texttt{<<<<<<< HEAD} y \texttt{======}.
    \item El desarrollador debe emplear los atajos visuales del IDE para seleccionar qué lógica impera (\textit{Accept Current Change} o \textit{Accept Incoming Change}).
    \item Purgados los marcadores, se debe persistir el cambio: \texttt{git add .}, confirmar (\texttt{git commit -m "fix: resuelve conflicto"}), y sincronizar (\texttt{git push origin tu-rama}). El PR se tornará verde automáticamente.
\end{enumerate}

\end{document}
